\section*{Rangkuman}
\addcontentsline{toc}{section}{Rangkuman}

Pada bab ini telah dibahas:
\begin{itemize}[leftmargin=*]
  \item pentingnya Teknik Kompilasi bagi sarjana informatika;
  \item posisi mata kuliah dalam kurikulum dan RPS berbasis OBE;
  \item arsitektur buku yang terdiri atas orientasi, landasan teori,
        empat belas bab inti, dan bab evaluasi;
  \item anatomi bab serta fungsi kotak pembelajaran dan kode program;
  \item proyek mini-kompilator Subset-C dan milestone-nya.
\end{itemize}

Bab berikutnya, Bab II, menyajikan landasan teori dan notasi formal
yang dipakai di seluruh buku: himpunan, relasi, bahasa, tata bahasa,
dan mesin abstrak. Materi tersebut menjadi bahasa bersama untuk
memahami analisis leksikal, sintaksis, dan semantik.
